% !PRINTED 590
\weberpage{590}
%<<CONT
ist die Partialnorm von $\mathfrak{p}$ eine Potenz von $\mathfrak{P}$. Wir setzen
\weq{3}{\mathfrak{N}_R(\mathfrak{p}) = \mathfrak{P}^f,}
und nennen $f$ \gesperrt{den Grad von $\mathfrak{p}$ in Bezug auf den Körper $R$.}

Die Norm von $\mathfrak{P}$ im Körper $R$ ist gleichfalls eine Potenz von $p$, die wir mit $P$ bezeichnen wollen, also
\weq{4}{N_R(\mathfrak{P}) = P.}

Es ergiebt sich dann für die Totalnorm von $\mathfrak{P}$ im Körper $\Omega$ nach \S.\,159, (5) % [sic] source names $\mathfrak P$ here, while (5) and the norm identity require $\mathfrak p$
\weq{5}{N_\Omega(\mathfrak{p}) = P^f.}

Ist $\mathfrak{p}^g$ die höchste in $\mathfrak{P}$ aufgehende Potenz von $\mathfrak{p}$, so ist $\mathfrak{P}$ nach 1. auch durch die $g^{\text{te}}$ Potenz aller mit $\mathfrak{p}$ conjugirten Primfactoren und durch keine höhere Potenz theilbar. Wenn nun $\mathfrak{p}_1, \mathfrak{p}_2, \ldots, \mathfrak{p}_e$ \gesperrt{die von einander verschiedenen} unter den mit $\mathfrak{p}$ conjugirten Primidealen sind, so ist
\weq{6}{\mathfrak{P} = (\mathfrak{p}_1\mathfrak{p}_2\ldots\mathfrak{p}_e)^g,}
und wenn man die Partialnorm auf beiden Seiten nimmt, und wie früher mit $n$ den Grad des Körpers $\Omega$ in Bezug auf $R$ bezeichnet, so dass die Partialnorm von $\mathfrak{P}$ gleich $\mathfrak{P}^n$ wird, so folgt aus (3):
\weq{7}{n = efg.}
Es sind also die natürlichen Zahlen $e, f, g$ Theiler von $n$.

Wenn $\mathfrak{p}$ durch $\varphi_1$ in $\mathfrak{p}_1$ übergeht, wenn also
\[
\mathfrak{p}_1 = \mathfrak{p} \,|\, \varphi_1,
\]
so ist $\mathfrak{p} = \mathfrak{p}_1 \,|\, \varphi_1^{-1}$, und wenn $\Psi$ die zu $\mathfrak{p}$ gehörige Gruppe ist, so ist auch für jede in $\Psi$ enthaltene Substitution $\psi$
\[
\mathfrak{p}_1 = \mathfrak{p} \,|\, \psi \varphi_1, \quad \mathfrak{p}_1 = \mathfrak{p}_1 \,|\, \varphi_1^{-1} \psi \varphi_1.
\]

Ist umgekehrt $\mathfrak{p}_1 \,|\, \varphi = \mathfrak{p}_1$, so folgt $\mathfrak{p} \,|\, \varphi_1 \varphi \varphi_1^{-1} = \mathfrak{p}$, d. h. $\varphi_1 \varphi \varphi_1^{-1} = \psi$ oder $\varphi = \varphi_1^{-1} \psi \varphi_1$.

Es gehört daher $\mathfrak{p}_1$ zur Gruppe $\varphi_1^{-1} \Psi \varphi_1$, und wenn
\[
\mathfrak{p}_1 = \mathfrak{p} \,|\, \varphi_1, \quad \mathfrak{p}_2 = \mathfrak{p} \,|\, \varphi_2, \ldots, \mathfrak{p}_e = \mathfrak{p} \,|\, \varphi_e
\]
ist, so ist
\weq{8}{\Phi = \Psi\varphi_1 + \Psi\varphi_2 + \cdots + \Psi\varphi_e.}
$\Psi$ ist also ein Theiler von $\Phi$ vom Index $e$ und vom Grade $gf$.
